\documentclass[11pt,a4paper]{article}
\usepackage[margin=1in]{geometry}
\usepackage[T1]{fontenc}
\usepackage{booktabs}
\usepackage{enumitem}
\usepackage{xcolor}
\usepackage{listings}
\usepackage[colorlinks=true,linkcolor=blue!60!black,urlcolor=blue!60!black]{hyperref}

\definecolor{codebg}{RGB}{245,247,247}
\definecolor{accent}{RGB}{78,141,126}

\lstset{
  basicstyle=\ttfamily\small,
  backgroundcolor=\color{codebg},
  frame=single,
  framerule=0pt,
  xleftmargin=8pt, xrightmargin=8pt,
  aboveskip=8pt, belowskip=8pt,
  breaklines=true,
  columns=fullflexible,
  keepspaces=true
}

\newcommand{\file}[1]{\texttt{#1}}

\title{\textbf{Personal Academic Website}\\[4pt]
\Large Maintenance, Updating \& GitHub Hosting Guide}
\author{Website of Md Zulquar Nain --- built with Quarto, R, and one Excel file}
\date{August 2026}

\begin{document}
\maketitle
\tableofcontents
\bigskip

\section{How the website is organised}

The project lives at \file{C:\string\Users\string\Abc\string\Desktop\string\academic-website}.
\textbf{Do not move it back into a Google Drive folder}: Drive's sync client
locks the \file{\_site} output folder during rendering and the render fails
with \emph{PermissionDenied} errors.

\begin{center}
\begin{tabular}{@{}ll@{}}
\toprule
\textbf{Path} & \textbf{What it is} \\
\midrule
\file{\_quarto.yml}            & Site configuration: navbar, footer, fonts, excluded folders \\
\file{index.qmd}               & Homepage (hero, interests, selected papers, data corner) \\
\file{research.qmd}            & Publications, projects, presentations (all from Excel) \\
\file{teaching.qmd}            & Teaching statement, course cards, courses taught, advising \\
\file{courses/}                & One \file{.qmd} per course website (hand-written, not Excel) \\
\file{Rclass\_2025/}           & Pre-built ECM2073 slide deck, served as-is at \file{/Rclass\_2025/} \\
\file{cv.qmd}                  & Web CV (awards, workshops, skills from Excel) \\
\file{blog.qmd} + \file{posts/} & Blog listing and the posts themselves \\
\file{R/helpers.R}             & \textbf{All} shared code; every page only calls functions from here \\
\file{R/pre-render.R}          & Runs before each render; refreshes \file{cv.pdf} from the CV project \\
\file{cvnain/}                 & The separate Quarto project that builds the PDF CV \\
\file{cvnain/data/cv\_data.xlsx} & \textbf{The single source of truth} (see below) \\
\file{background.js}           & The animated homepage background \\
\file{irf-surface.js}          & Generated: the estimated IRF surface it draws (Section~\ref{sec:irfdata}) \\
\file{styles.css}              & Palette, fonts, cards, badges, reading measure \\
\file{images/}                 & Profile photo and other pictures \\
\file{archive/}                & Not needed to build or host the site; kept, never committed \\
\file{.github/workflows/publish.yml} & Auto-deployment to GitHub Pages \\
\file{.github/workflows/check-inflation-data.yml} & Monthly nag when the inflation tracker falls behind \\
\file{software/}               & Shiny splicer app and the inflation tracker \\
\file{\_site/}                 & Rendered output --- never edit by hand \\
\bottomrule
\end{tabular}
\end{center}

\section{Updating content: one Excel file}

\file{cvnain/data/cv\_data.xlsx} drives \emph{both} the PDF CV and the
website. Each sheet feeds specific parts of both documents:

\begin{center}
\begin{tabular}{@{}lll@{}}
\toprule
\textbf{Sheet} & \textbf{Feeds on the website} & \textbf{Feeds in the PDF CV} \\
\midrule
publications     & Research page + homepage cards & Publications section \\
presentations    & Research page                  & Presentations / Conferences \\
projects         & Research page                  & Projects \\
teaching\_courses & Teaching page (by level)      & Teaching section (by level) \\
advising         & Teaching page                  & Advising \\
awards           & CV page                        & Scholarship / Awards \\
workshops        & CV page                        & Workshops and Training \\
service          & CV page                        & (currently CV-off) \\
memberships      & CV page                        & Academic Affiliation \\
skills           & CV page                        & Skills \\
scholar\_fallback & Stats badges (fallback tier)  & Stats box (fallback tier) \\
software         & Software page cards (incl.\ \texttt{status\_file}) & --- \\
\bottomrule
\end{tabular}
\end{center}

\subsection{Typical tasks}

\textbf{Add a new paper.} Add a row to the \emph{publications} sheet
(\texttt{category} = \texttt{journal}, \texttt{submitted}, or
\texttt{working}; fill \texttt{authors}, \texttt{year}, \texttt{title},
\texttt{journal}, \texttt{doi}). Re-render (Section~\ref{sec:render}).

\textbf{Choose the papers featured on the homepage.} Add a column named
\texttt{selected} to the publications sheet and put \texttt{YES} on up to
three rows. Without it, the three most recent journal articles are shown.

\textbf{Update Scholar fallback figures.} Edit the \emph{scholar\_fallback}
sheet occasionally (citations, h\_index, i10\_index, updated). It is used
whenever Google Scholar blocks automated access (which is common --- see
Section~\ref{sec:scholar}).

\textbf{Things not in Excel} (they rarely change): hero text and interest
cards in \file{index.qmd}; education and appointments in \file{cv.qmd};
the teaching statement in \file{teaching.qmd}; everything on the course
pages in \file{courses/} (Section~\ref{sec:courses}).

\textbf{Why course pages are not in Excel.} Excel earns its place when the
same rows feed \emph{two} outputs --- publications and teaching\_courses
print into both the website and the PDF CV, so one edit updates both. Course
material feeds one page only, and a course page is mostly prose (syllabus,
prerequisites, assessment) that no spreadsheet row can hold. So courses are
plain \file{.qmd} files.

\subsection{Writing a blog post}

Create a folder under \file{posts/} with an \file{index.qmd}:

\begin{lstlisting}
posts/my-new-post/index.qmd
---
title: "My new post"
description: "One-line teaser shown in the listing."
date: 2026-08-01
categories: [economics, teaching]
---
Body of the post in ordinary markdown...
\end{lstlisting}

It appears on the Blog page automatically at the next render.

\section{Course pages}\label{sec:courses}

\subsection{How the Teaching page is put together}

\file{teaching.qmd} has four parts, top to bottom:

\begin{enumerate}[itemsep=1pt]
  \item The \textbf{teaching statement} --- ordinary prose at the top of the
        file. Edit it directly.
  \item \textbf{Course Websites} --- a card for every \file{.qmd} in
        \file{courses/}. This list is \emph{generated}; there is nothing to
        maintain.
  \item \textbf{Courses Taught} tables --- from the \emph{teaching\_courses}
        Excel sheet, unchanged.
  \item \textbf{Advising} --- from the \emph{advising} Excel sheet, unchanged.
\end{enumerate}

The cards come from the \texttt{listing:} block in the front matter of
\file{teaching.qmd}, the same mechanism the Blog page uses for
\file{posts/}. It scans \file{courses/} and builds one card per file, newest
\texttt{date} first.

\subsection{Adding a new course}\label{sec:newcourse}

Create one file, \file{courses/my-course.qmd}. Everything in the front
matter feeds the card on the Teaching page:

\begin{lstlisting}
---
title: "Financial Econometrics"
subtitle: "ECM3081 . M.A. Economics . Aligarh Muslim University"
description: "One or two lines describing the course. This is the
  paragraph shown on the Teaching page card."
date: 2026-01-10          # newest date sorts to the top
image: /images/course-fin-econ.jpg
categories: [R, Econometrics, PG]
toc: true
---

## About the course
Ordinary markdown from here down: syllabus, readings, links...
\end{lstlisting}

Render, and the Teaching page card appears on its own. \textbf{There is no
list to update anywhere} --- which is the point: the index can never fall
out of step with the pages that actually exist.

Put the card image in \file{images/} and refer to it as
\file{/images/name.jpg} (the leading slash means ``from the project root'').
Roughly 800$\times$500 pixels works well. If you omit \texttt{image:}
entirely the card still renders, just without a thumbnail.

\subsection{Linking to slides and data}\label{sec:slidelinks}

Course material lives in its own folder at the top level of the project
(\file{Rclass\_2025/} is the first one). Because the course page sits one
level down in \file{courses/}, links start with \texttt{../}:

\begin{lstlisting}
[Open](<../Rclass_2025/25_Lecture 1 Introduction to R.html>)
\end{lstlisting}

\textbf{The angle brackets matter.} Your slide filenames contain spaces.
Ordinary markdown link syntax breaks on a space; wrapping the path in
\texttt{<...>} tells Pandoc to take the whole thing as one URL, so you can
paste the filename exactly as it appears in Windows Explorer and never
think about \texttt{\%20}.

\textbf{Capitalisation matters too.} Windows treats
\file{25\_lecture 1...} and \file{25\_Lecture 1...} as the same file;
the Linux server behind GitHub Pages does not. A link that works locally
can 404 once published. Copy filenames, do not retype them.

\subsection{Publishing a folder of pre-built slides}

\file{Rclass\_2025/} was rendered in its own RStudio project and its HTML is
finished. Two entries in \file{\_quarto.yml} publish it untouched:

\begin{lstlisting}
render:
  - "!Rclass_2025/"          # do NOT re-render the .qmd/.Rmd inside
resources:
  - Rclass_2025/**           # DO copy the finished HTML and assets
  - "!Rclass_2025/_site/**"  # skip the duplicate copy of the same deck
\end{lstlisting}

The \texttt{render} exclusion is essential: without it Quarto tries to
rebuild all sixteen slide decks on every site render, which is slow and
fails on the two \texttt{xaringan} decks. To add a second course folder
later, copy both lines and change the name.

\subsection{Rebuilding a slide deck}\label{sec:xaringan}

Most of the ECM2073 decks are Quarto \texttt{revealjs}. Two of them
(Lecture 8 \emph{R Base Plot} and Lecture 12 \emph{ggplot2}) are older
\texttt{xaringan} decks written in R Markdown, and they are built by
\textbf{rmarkdown}, not Quarto. From the \file{Rclass\_2025} folder:

\begin{lstlisting}
$env:RSTUDIO_PANDOC = "C:\Program Files\RStudio\resources\app\bin\quarto\bin\tools"
& "C:\Program Files\R\R-4.5.2\bin\Rscript.exe" -e "rmarkdown::render('25_Lecture 8 R Base Plot.Rmd')"
\end{lstlisting}

The \texttt{RSTUDIO\_PANDOC} line is the part people miss: R on its own
cannot find Pandoc outside RStudio and fails with \emph{``pandoc version
1.12.3 or higher is required''}. Inside RStudio just press \textbf{Knit}
and none of this applies.

These decks need \texttt{xaringan}, \texttt{tidyverse}, \texttt{skimr},
\texttt{janitor}, \texttt{DT}, \texttt{here}, \texttt{kableExtra} and
\texttt{titanic}. They write a shared \file{Rclass\_2025/libs/} folder plus a
\file{..\_files/} folder of figures --- both must be committed alongside the
HTML or the published slides lose their styling and images.

\subsection{Size warning}\label{sec:size}

\file{Rclass\_2025/} used to be dominated by two screencasts. Both were
recorded at a fixed high bitrate (17.6 and 8.3 Mbit/s) on static screen
content that needs a small fraction of it, and both have been re-encoded
in place at \texttt{crf 28} with no change of resolution, frame rate or
duration:

\begin{center}
\begin{tabular}{@{}lrr@{}}
\toprule
\textbf{File} & \textbf{Was} & \textbf{Now} \\
\midrule
\file{images/dataimport.mp4} (Lecture 5) & 86 MB & 0.9 MB \\
\file{images/project1.mp4} (Lecture 2)   & 50 MB & 1.1 MB \\
Everything else                          & \multicolumn{2}{r}{$\approx$108 MB} \\
\midrule
\textbf{Folder total}                    & 244 MB & $\approx$110 MB \\
\bottomrule
\end{tabular}
\end{center}

Quality was checked before the swap: SSIM 0.998 and PSNR 46 dB against the
originals, and RStudio console text is indistinguishable at 1:1. The
originals were moved (not deleted) to
\file{academic-website-video-originals/} next to the project folder.

The whole published site now comes to about \textbf{233 MB}
(\file{Rclass\_2025/} 110 MB, the Shinylive splicer app 96 MB, the
inflation tracker's libraries 21 MB, everything else 6 MB), against the
\textbf{1 GB} that GitHub Pages asks a published site to stay under.

What the \emph{repository} carries is smaller again. Git stores blobs by
content, so the sixteen \file{*\_files/} folders in \file{Rclass\_2025/} ---
1408 files, but only 290 distinct ones, because every deck ships its own
copy of reveal.js, Font Awesome and Source Sans Pro --- cost about 30 MB
once, not 110 MB. There is no reason to consolidate them by hand.

Across the whole project the tracked working tree is about 145 MB but only
53 MB of it is distinct content, and that --- not the 145 --- is what the
first push actually carries.

No file in the project now exceeds 50 MB, the threshold at which GitHub
starts warning; the hard rejection limit is 100 MB. The 247 MB duplicate
render of the deck has been moved to \file{archive/Rclass\_2025/\_site/} and
can be deleted from disk entirely whenever you like.

\section{Rendering}\label{sec:render}

\subsection{From RStudio (easiest)}
Open \file{academic-website} as the working directory and run
\texttt{quarto render} in the Terminal pane (or use Build $\to$ Render
Website). RStudio finds R automatically.

\subsection{From a plain terminal}
The RStudio-bundled Quarto needs to be told where R lives:

\begin{lstlisting}
$env:QUARTO_R = "C:\Program Files\R\R-4.5.2\bin"
& "C:\Program Files\RStudio\resources\app\bin\quarto\bin\quarto.exe" render
\end{lstlisting}

Use \texttt{quarto preview} instead of \texttt{render} for a live-reload
local server while editing.

\subsection{Updating the PDF CV}
Render the CV inside its own project first, then the site:

\begin{lstlisting}
cd cvnain
quarto render "Dr Nain_CV.qmd"     # produces Dr-Nain_CV.pdf
cd ..
quarto render                       # pre-render copies it to cv.pdf
\end{lstlisting}

\section{Google Scholar, OpenAlex, and caching}\label{sec:scholar}

The stats badges (Citations, h-index, i10-index) are resolved in this
order at render time:

\begin{enumerate}[itemsep=1pt]
  \item \textbf{Live scrape} of the Scholar profile
        (ID \texttt{7Db2EHMAAAAJ}). Google frequently answers scripted
        requests with a CAPTCHA, in which case this tier fails silently.
  \item \textbf{Cache} of the last successful live fetch
        (\file{data/scholar\_stats\_cache.csv}) --- written automatically
        whenever tier 1 succeeds. \emph{Commit these files to git} so the
        GitHub build (whose datacenter IPs are always blocked) can use them.
  \item \textbf{Excel} \emph{scholar\_fallback} sheet.
  \item \textbf{OpenAlex} (author ID \texttt{A5034570980}), a free API that
        is never blocked --- used only if the Excel values are missing.
        Note its counts are lower than Scholar's (different coverage).
\end{enumerate}

\section{The Data Corner (MoSPI)}\label{sec:mospi}

The homepage chart shows \textbf{monthly headline CPI inflation} and
\textbf{IIP growth} for All-India, fetched at render time from the Ministry
of Statistics and Programme Implementation's open API
(\texttt{api.mospi.gov.in}) --- the same backend as the
\href{https://esankhyiki.mospi.gov.in}{eSankhyiki portal}, and the same one
\file{software/inflation-tracker.qmd} uses. \textbf{No API key is needed.}
Each successful fetch caches to \file{data/mospi\_cache.csv}, which is used
if the API is unreachable, so commit that file.

It replaced an annual World Bank series in August 2026: World Bank data for
India lags by a year or more, while MoSPI publishes CPI around the 12th of
each month. \file{data/wdi\_cache.csv} is left over from that and can be
deleted.

\subsection{Knobs}

\begin{center}
\begin{tabular}{@{}lp{8.5cm}@{}}
\toprule
\textbf{In \file{R/helpers.R}} & \textbf{Effect} \\
\midrule
\texttt{ECON\_YEARS} & How many years of history the chart shows (currently 5) \\
\texttt{cpi\_series()} & Which CPI cut is plotted --- currently All-India, Combined sector, headline \\
\texttt{iip\_series()} & Currently the General index; \texttt{type} can also be \emph{Sectoral} or \emph{Use-based category} \\
\texttt{smoothing} in \texttt{econ\_chart()} & Curve of the line: 0 straight, 1 current, 1.3 maximum \\
\bottomrule
\end{tabular}
\end{center}

\subsection{Why the series are chain-linked}\label{sec:chain}

MoSPI republishes each index on a new base every few years, and the old and
new series \emph{disagree where they overlap}. The base 2022--23 IIP puts
March 2024 growth at \textbf{19.0\%}; base 2011--12 puts the same month near
\textbf{5\%}, because the new base's first year is measured against a partial
history. Taking ``the newest base wins'' would draw that 19\% spike as though
it were real.

\texttt{chain\_bases()} therefore takes the \emph{oldest} base first and lets
each newer base contribute only the months the previous one does not reach:
the long-running series carries the history, the newest one extends the tip,
and the join sits at the end where the discontinuity is smallest. If you add
another indicator, route it through the same function.

\textbf{Not every spike is an artifact.} The chart's May 2022 IIP reading of
19.7\% is genuine --- a base effect against the 2021 Delta-wave lockdown ---
and is left in.

\subsection{Other Indian sources}

If you want to add series later, in rough order of how pleasant they are to
work with:

\begin{center}
\begin{tabular}{@{}lp{7.4cm}l@{}}
\toprule
\textbf{Source} & \textbf{Covers} & \textbf{Key?} \\
\midrule
MoSPI API      & CPI, IIP (1994--), ASI; monthly & no \\
data.gov.in    & Very broad catalogue; quality varies a lot by dataset & free key \\
RBI DBIE       & Monetary, banking, BoP, fiscal --- the richest, but no documented public API; scraping breaks on redesign & --- \\
World Bank     & Cross-country, annual, lagged & no \\
UN Comtrade    & Bilateral trade (\texttt{comtradr}) & free key \\
BIS            & Policy rates, REER, credit-to-GDP & no \\
\bottomrule
\end{tabular}
\end{center}

\section{The inflation tracker}\label{sec:tracker}

\file{software/inflation-tracker.qmd} is a standing dashboard of Indian
retail inflation, published on the Software page as the second card. It is
\textbf{rendered by hand}, not by the site build.

\subsection{Why it is not automated}

It needs \texttt{sf} (and therefore GDAL, GEOS and PROJ as system
libraries), plus \texttt{ggiraph}, \texttt{ragg} and \texttt{DT}. Installing
that on a GitHub runner turns a four-minute build into fifteen, for a page
that changes \emph{twelve times a year}. More to the point, MoSPI's API is
young --- base years and field names shift --- and an automated publish
fails \emph{silently}: the build goes green and a broken page sits on your
site. Rendering by hand keeps you looking at the state map at the moment
things break.

\subsection{The monthly routine}

MoSPI publishes CPI around the 12th at 16:00 IST. After that:

\begin{lstlisting}
cd software
quarto render inflation-tracker.qmd
cd ..
quarto render
git add . ; git commit -m "Inflation tracker: <month> data" ; git push
\end{lstlisting}

Look at the state map before pushing. It is the first thing to break when
a field name changes.

\subsection{Three parts that must stay in step}

\begin{center}
\begin{tabular}{@{}lp{8.2cm}@{}}
\toprule
\textbf{Piece} & \textbf{Role} \\
\midrule
\file{software/tracker\_status.json} & Written on every tracker render:
  latest data month, headline rate, render date \\
\texttt{print\_software()} & Reads it via the \texttt{status\_file} column
  in the \emph{software} sheet and stamps ``Data through Jul 2026'' on the
  card, so a stale page announces itself \\
\file{.github/workflows/} \file{check-inflation-data.yml} & Compares it
  against MoSPI daily on the 12th--18th and opens an issue when the page
  is behind \\
\bottomrule
\end{tabular}
\end{center}

The status file is the hinge: rename a field in it and both the card stamp
and the notifier stop working. It must be committed.

\subsection{The notifier}

The Action \emph{checks but never publishes}. It runs at 12:00 IST daily
from the 12th to the 18th, asks MoSPI for the newest CPI month carrying an
inflation figure, and compares it with \file{tracker\_status.json}. If the
site is behind it opens a labelled issue with the rebuild commands; it will
not open a second issue for the same release. Nothing else in the repository
depends on it, so if it ever becomes a nuisance, delete the file.

Run it on demand from the repository's \emph{Actions} tab
(\emph{Run workflow}) to test it.

\subsection{Why the page is not self-contained}

\texttt{embed-resources} is deliberately \texttt{false}. Self-contained, the
tracker is one 11 MB file, and a \emph{fresh} 11 MB blob enters git history
at every monthly rebuild --- about 130 MB a year, permanently. Split out,
the DataTables and \texttt{ggiraph} JavaScript in
\file{inflation-tracker\_files/libs/} (20 MB) is byte-identical each render
and git stores it once; only the six charts churn, about 1 MB a month.
The cost is a one-off 20 MB and the rule that
\file{inflation-tracker\_files/}, \file{tracker.css} and
\file{tracker\_status.json} must all be published alongside the HTML ---
they are listed in \file{\_quarto.yml} under \texttt{resources}.

\bigskip

Two MoSPI endpoints exist but were not usable as of August 2026:
\texttt{/api/nas/getNASData} (national accounts) returns HTTP 500 for every
parameter combination tried, and \texttt{/api/asi/getASIData} wants a
\texttt{classification\_year} whose accepted values are undocumented. Worth
re-probing later --- the API tells you what it wants if you call it with no
parameters and read the error body.

\section{Customising the look}

\subsection{The animated background (\file{background.js})}
Four layers share one full-page canvas: a drifting network of variable
symbols, floating equations, small 2-D impulse-response panels that draw
themselves and respawn, and the 3-D surface. Everything is controlled by the
\texttt{STYLE} block at the top of the file:

\begin{center}
\begin{tabular}{@{}lp{8.4cm}@{}}
\toprule
\textbf{Knob} & \textbf{Effect} \\
\midrule
\texttt{intensity} & Master visibility: 0.5 subtle, 1 current, 2 bold \\
\texttt{speed}     & Master animation pace \\
\texttt{palette}   & Five [R,G,B] colours used by equations, nodes, panels \\
\texttt{surface.size / posX / posY} & Size and position of the 3-D surface \\
\texttt{surface.faceAlpha} & Opacity of the shaded faces. Raise for a more
  solid surface, lower if it competes with the hero text \\
\texttt{surface.edgeAlpha} & Opacity of the mesh drawn over the faces \\
\texttt{surface.colormap} & \texttt{"parula"} (MATLAB's default) or
  \texttt{"site"} to use the page palette instead \\
\texttt{surface.azimuth / elevation} & The fixed viewpoint. Currently
  MATLAB's \texttt{view(-46, 29)}. There is no spin \\
\texttt{surface.evolve} & How fast the \emph{fallback} surface reshapes;
  unused when real data is present \\
\texttt{EQUATIONS / LABELS} arrays & The floating formulas and node symbols \\
\bottomrule
\end{tabular}
\end{center}

The surface is drawn like MATLAB's \texttt{surf()}: faces filled from the
colormap with a mesh over them, painted back to front. Colour encodes the
size of the response and is rescaled to the data's own range each frame,
which is what MATLAB's automatic \texttt{caxis} does --- against a fixed
ceiling the surface only ever samples the blue end and looks washed out.

\paragraph{Do not add per-frame motion.} Every speed in the file is
multiplied by \texttt{step}, the number of 60\,Hz frames the last frame
actually took. Without it a 120\,Hz monitor runs the animation at double
speed and 144\,Hz at 2.4$\times$. If you add anything that moves, multiply
its increment by \texttt{step} as well.

The animation also pauses when the tab is hidden \emph{and} when the canvas
scrolls out of view, and it removes itself if the reader has asked for
reduced motion.

\subsection{The estimated IRF surface (\file{irf-surface.js})}
\label{sec:irfdata}

The 3-D surface is not decoration: it is the estimated response of inflation
to an activity shock, from the hybrid TVP-BVAR with stochastic volatility in
\file{Desktop/philips}. The numbers live in \file{irf-surface.js}, a
generated file of about 5\,KB that assigns one global,
\texttt{window.\_\_IRF\_SURFACE\_\_}, and is loaded by \file{index.qmd}
\emph{before} \file{background.js}.

It is delivered as a script rather than JSON on purpose: \texttt{fetch()} is
blocked under \texttt{file://}, so a JSON file would silently fall back to
the invented surface whenever \file{\_site/index.html} is opened from disk.
If the file is missing altogether the homepage still works --- it renders a
parametric stand-in instead.

\paragraph{Regenerating it.} Only needed if the model is re-estimated.

\begin{enumerate}[itemsep=2pt]
\item In MATLAB, load the chosen results file, redo STEP 7 of
      \file{run\_phillips\_curve.m} and save the array:
\begin{lstlisting}
load('India/Results_modspec_74.mat');
[theta, hret] = shuffleparameters(store_beta, store_alp, ...
                                  store_h, included_vars, p);
htmp = hret; model = 1;   % is_sv = [1 1] -> stochastic volatility
impulses = impulseresp(model, n, p, theta, [], htmp, 20, 1:T, 1, ...
                       0, 0, 0, 0, 0, OnlyStationary);
save('irf_surface.mat', 'impulses', '-v7.3');
\end{lstlisting}
      This takes about 75 seconds. The MCMC itself is \emph{not} re-run.
\item In R, read that file and write the script. \texttt{impulses} is
      indexed (response, shock, horizon, time); response 2 is inflation and
      shock 1 is activity, because the runner assembles
      \texttt{data = [act\_col, inf\_data]}.
\end{enumerate}

\paragraph{Three traps, all of which cost time once already.}
\begin{itemize}[itemsep=2pt]
\item \texttt{comp\_resp = 0} in \file{user\_config.m} means the runner never
      computes IRFs at all. That is why no impulse responses are stored in
      the \file{Results\_modspec\_*.mat} files --- they hold only the MCMC
      draws.
\item Passing \texttt{ci = 0} returns a \emph{4-D} array, not the 5-D one
      that \file{make\_figures.m}'s \texttt{impulses(i,j,:,:,1)} implies.
      Ask for credible intervals and the quantile dimension reappears.
\item MAT v7.3 is HDF5, and the usual advice is that HDF5 stores dimensions
      reversed. \textbf{Do not transpose.} \texttt{hdf5r} already undoes the
      reversal on read: its dims match MATLAB's \texttt{size()} exactly, and
      applying \texttt{aperm()} silently turns the $21\times28$ surface into
      a $2\times2$ slice.
\item \texttt{OnlyStationary = 1} discards roughly 5{,}300 of the 20{,}000
      draws at each time point, so the surface is a median over about
      14{,}700 draws. This is expected, not an error.
\end{itemize}

Heights are divided by the largest absolute response rather than shifted, so
zero stays at zero and the genuinely negative part of the response dips
below the base plane instead of being flattened onto it. Because estimated
data does not evolve, the surface sweeps itself in once along the time axis
and then holds.

\subsection{Site styling (\file{styles.css})}
The \texttt{.section} background alpha (currently 0.80) controls how much
of the animation shows through content sections. Scholar badges are styled
by \texttt{.sbadge}.

\subsubsection*{The palette has three tiers --- pick by contrast}

The teal comes in three values, and \textbf{which one you use depends on
how large the text is}, not on taste:

\begin{center}
\begin{tabular}{@{}llp{6.4cm}@{}}
\toprule
\textbf{Token} & \textbf{On white} & \textbf{Use for} \\
\midrule
\texttt{--accent}      & 2.5:1 & Fills, rules, borders. \textbf{Never text.} \\
\texttt{--accent-dark} & 3.9:1 & Large text, borders \\
\texttt{--accent-text} & 6.2:1 & Small text and links \\
\texttt{--ink}         & 14.7:1 & Body text \\
\texttt{--ink-soft}    & 7.6:1 & Secondary text \\
\bottomrule
\end{tabular}
\end{center}

WCAG AA wants \textbf{4.5:1} for normal text. \texttt{--accent-dark} does
not reach it, which is why \texttt{--accent-text} was added in August 2026
--- \texttt{.paper-meta} had been setting 0.84\,rem uppercase text at
3.9:1, the hardest case to read. If you introduce a new colour, check it
before using it for text; there are contrast checkers online, or in R:

\begin{lstlisting}
lum <- function(hex) { v <- col2rgb(hex)[,1]/255
  f <- ifelse(v <= 0.03928, v/12.92, ((v+0.055)/1.055)^2.4)
  sum(c(0.2126,0.7152,0.0722)*f) }
(lum("#ffffff")+0.05) / (lum("#376a5d")+0.05)   # 6.21
\end{lstlisting}

\subsubsection*{Overriding the Quarto theme}

The \texttt{cosmo} theme ships a blue link colour (\texttt{\#2761e3}), a
blue callout accent and pink inline code --- three accents fighting the
teal. \file{styles.css} pulls all three onto the palette under
\texttt{main a}, \texttt{main .callout-note} and
\texttt{main code:not(.sourceCode)}. Buttons and the navbar set their own
colours with \texttt{!important} and are unaffected.

The callout icon is a blue SVG recoloured with \texttt{hue-rotate}. Blue
sits near 220\textdegree; \texttt{-60deg} lands on the teal. Rotating the
other way turns it pink.

\subsubsection*{Reading measure}

\begin{lstlisting}
.page-layout-article main.content { max-width: 72ch; line-height: 1.7; }
\end{lstlisting}

Prose past roughly 72 characters a line is tiring to read, and Bootstrap's
1.5 line-height is tight for long citation lists. \textbf{The
\texttt{.page-layout-article} prefix is essential}: the homepage carries
\texttt{page-layout-custom} and its full-bleed sections must not be
constrained. Check which class a page has by searching its rendered HTML
for \texttt{id="quarto-content"}.

\subsubsection*{Two gotchas worth remembering}

\begin{description}[style=nextline,itemsep=3pt]
\item[Listing ids are renamed]
  Quarto turns \texttt{::: \{\#course-pages\}} into a div with id
  \texttt{listing-course-pages}. Styling \texttt{\#course-pages} silently
  matches nothing --- which is exactly how the Teaching page card went
  unstyled for a while. Style \texttt{\#listing-course-pages}.
\item[\texttt{.paper-meta} is uppercase]
  Fine for a short label such as the tech line on a software card;
  unreadable for a whole sentence. Sentences use \texttt{.soft-note}.
\end{description}

\subsubsection*{The publication list}

The Research page is not a numbered list --- it is a two-column grid built
by \texttt{pub\_entry\_html()} in \file{R/helpers.R}: year in the left
gutter, then title, authors, and venue/DOI in quieter type. The year prints
only on the \emph{first} paper of each year so runs group visually. Styling
lives under \texttt{.pub-list}, \texttt{.pub-entry}, \texttt{.pub-year},
\texttt{.pub-title}, \texttt{.pub-authors} and \texttt{.pub-venue}. Titles
are styled on \texttt{.pub-title, .pub-title a} together so a paper with a
DOI looks identical to one without.

\subsection{Checking how a change looks}\label{sec:screenshot}

\texttt{quarto preview} gives a live local server. To capture pages as
images without opening a browser --- useful for comparing before and after,
or for checking a page you are not currently looking at:

\begin{lstlisting}
$chrome = "C:\Program Files\Google\Chrome\Application\chrome.exe"
& $chrome --headless=new --disable-gpu --hide-scrollbars `
  --window-size=1440,1400 --virtual-time-budget=8000 `
  --screenshot="shot.png" `
  "file:///C:/Users/Abc/Desktop/academic-website/_site/research.html"
\end{lstlisting}

\texttt{--virtual-time-budget} gives the page time to run its JavaScript,
which the homepage's animated background and the plotly charts need.
Note that over \texttt{file://} the browser blocks remote images, so the
footer's visitor-count badge shows as broken in these screenshots --- that
is an artifact of local capture, not a fault on the live site.

\subsection{Section formatting (\file{R/helpers.R})}
How each Excel-driven section is printed lives here: publication section
headings in \texttt{PUB\_SECTIONS}, course level headings in
\texttt{TEACH\_LEVELS}, Data Corner settings in \texttt{ECON\_YEARS} and
\texttt{MOSPI\_BASE} (Section~\ref{sec:mospi}), chart colours in
\texttt{CHART\_COLORS}.

\subsection{Courses taught}
Courses are grouped by level (\emph{Undergraduate}, \emph{Postgraduate},
\emph{Ph.D.}) and run together into one semicolon-separated sentence each,
rather than set as a table --- the old table spent a whole column repeating
three values down twelve rows. To add a course, add a row to
\emph{teaching\_courses} with \texttt{institution}, \texttt{course\_title}
and \texttt{level}; it is filed under the right heading automatically and
sorted alphabetically. A \texttt{level} code that is not \texttt{UG},
\texttt{PG} or \texttt{PhD} still prints, under its own raw code, so a typo
shows up on the page rather than vanishing.

The same grouping is produced twice, by two pieces of code that must be
kept in step: \texttt{TEACH\_LEVELS} plus \texttt{print\_teaching()} in
\file{R/helpers.R} for the website, and the Teaching chunk inside
\file{cvnain/Dr Nain\_CV.qmd} for the PDF. If you rename a heading, rename
it in both.

\section{Hosting on GitHub Pages}

The GitHub account is \url{https://github.com/alignain}.

\subsection{Choosing the repository name --- this fixes the URL}

Two options, and they give different addresses:

\begin{center}
\begin{tabular}{@{}llp{5.4cm}@{}}
\toprule
\textbf{Repository} & \textbf{Site address} & \textbf{Notes} \\
\midrule
\texttt{alignain.github.io} & \texttt{https://alignain.github.io/} &
  The clean, root-level address. One per account. \textbf{Recommended.} \\
\texttt{website} (any name) & \texttt{https://alignain.github.io/website/} &
  Everything sits under a sub-path. \\
\bottomrule
\end{tabular}
\end{center}

Pick before the first push: changing later means every link anyone has
saved breaks. The rest of this section assumes \texttt{alignain.github.io}.

\subsection{First-time setup}

\paragraph{Step 0 --- install git.} It is not currently on this machine ---
rechecked August 2026: no \texttt{git}, no \texttt{gh}, and nothing bundled
inside RStudio. Either
\begin{lstlisting}
winget install --id Git.Git -e
\end{lstlisting}
or take the installer from \url{https://git-scm.com/download/win} and accept
the defaults. Reopen the terminal afterwards so \texttt{git} is on the
\texttt{PATH}. Then set your identity once:
\begin{lstlisting}
git config --global user.name  "Md Zulquar Nain"
git config --global user.email "alignain@gmail.com"
\end{lstlisting}

\paragraph{Step 1 --- create the repository.} On
\url{https://github.com/new}, owner \texttt{alignain}, name
\texttt{alignain.github.io}, visibility \textbf{Public} (Pages needs a paid
plan to publish from a private repository). Add \emph{no} README,
\emph{no} \file{.gitignore} and \emph{no} licence --- the folder already has
what it needs, and an initial commit on GitHub's side only causes a
conflict on the first push.

\paragraph{Step 2 --- push the folder.} From \file{academic-website}:
\begin{lstlisting}
git init
git add .
git status --short | measure          # sanity check: how many files?
git commit -m "Academic website"
git branch -M main
git remote add origin https://github.com/alignain/alignain.github.io.git
git push -u origin main
\end{lstlisting}
The first push carries roughly 50\,MB and will take a few minutes. Git will
ask you to sign in to GitHub in a browser window; that is normal, and it
stores the credential so later pushes are silent.

\paragraph{Step 3 --- turn on Pages.} On GitHub:
\emph{Settings} $\to$ \emph{Pages} $\to$ \emph{Build and deployment} $\to$
\emph{Source} $\to$ \textbf{GitHub Actions}. \textbf{Nothing is published
until this is set}, and the deploy step of the workflow will fail with a
permissions error if it is left on the default.

There is no \texttt{gh-pages} branch. The workflow renders the site, uploads
\file{\_site} as a Pages artifact and deploys from that, so the built output
never enters the repository history. Do \emph{not} run
\texttt{quarto publish gh-pages}: it would create a branch that competes
with the Actions deployment.

\paragraph{Step 4 --- watch the first build.} The \emph{Actions} tab shows
\emph{Publish website to GitHub Pages} running. Expect \textbf{10--20
minutes} the first time: it installs R and twelve packages, then downloads
about 115\,MB of WebAssembly runtime to rebuild the Shinylive app. Later
builds reuse a cache and take a few minutes.

\paragraph{Step 5 --- the address is already recorded.} This step used to
say ``set \texttt{site-url} once the run is green''. It is now set in
\file{\_quarto.yml} ahead of the first push:
\begin{lstlisting}
site-url: https://alignain.github.io/
\end{lstlisting}
so there is nothing to do here and no second push. Quarto uses it for
absolute links and social previews, and it is what makes Quarto emit
\file{\_site/sitemap.xml} and \file{\_site/robots.txt} at all --- with
\texttt{site-url} unset, neither file is written and search engines get no
index. If you ever change the repository name, change this line to match, or
every address in the sitemap will point somewhere that does not exist.

\subsection{Every update afterwards}

\begin{lstlisting}
git add .
git commit -m "Update publications"
git push
\end{lstlisting}

The included GitHub Action (\file{.github/workflows/publish.yml}) installs
R, the required packages (\texttt{readxl, dplyr, glue, plotly, httr2,
jsonlite, rvest, httr, tidyr, knitr, rmarkdown, shinylive}), and Quarto,
rebuilds the Shinylive app, re-renders the site in the cloud, and deploys
it. The first cloud build takes several minutes --- it also downloads
roughly 115\,MB of WebAssembly runtime for the Shinylive app --- while
later builds reuse a cache and are faster.

\subsection{What to commit}
Commit everything except what \file{.gitignore} already excludes
(\file{\_site/}, \file{.quarto/}, R session files, \file{archive/},
\file{.claude/} --- per-machine editor settings, not project content --- and
\file{software/splicer-app/} --- 115\,MB of generated WebAssembly that the
cloud build regenerates from \file{software/app.R} with
\file{software/build-shinylive.R}).

This guide is the one exception worth spelling out, because it is the only
hand-run \LaTeX{} project in the repository. Its \file{.tex} source
\emph{and} the built \file{Website\_Maintenance\_Guide.pdf} are both
committed --- the PDF is the copy anyone actually opens, and no cloud build
compiles it --- while pdflatex's scratch files
(\file{.aux}, \file{.log}, \file{.out}, \file{.toc}, and the
\file{.fls}/\file{.fdb\_latexmk}/\file{.synctex.gz} that latexmk and editor
integrations add) are ignored. Those are rewritten on every compile and mean
nothing to anyone else. The rule is written as \file{/guide/*.aux} and so on
rather than a bare \file{*.log}, which would quietly swallow unrelated logs
elsewhere in the project.

After editing the guide, recompile it \emph{twice} --- the second pass is
what resolves the table of contents and the \texttt{Section~n} cross-references --- and commit the new
PDF alongside the source:
\begin{lstlisting}
cd guide
pdflatex -interaction=nonstopmode Website_Maintenance_Guide.tex
pdflatex -interaction=nonstopmode Website_Maintenance_Guide.tex
\end{lstlisting}
pdflatex is MiKTeX's, installed under
\file{\%LOCALAPPDATA\%\string\Programs\string\MiKTeX}.

\file{archive/} holds everything not needed to build or host the site: the
legacy \file{cv/} project, the duplicate deck render, and R session junk.
Nothing there was deleted, and \file{archive/README.md} says what each item
is and how to put it back.

In particular \emph{do} commit \file{irf-surface.js} (the estimated surface
--- it is generated, but from a 75-second MATLAB step that the cloud build
cannot run), \file{data/*.csv} caches and
\file{cvnain/} (the CV project and its Excel file are needed by the cloud
build), and \file{Rclass\_2025/} minus its own \file{\_site/} sub-folder
--- the cloud build copies those slides straight into the published site,
so if they are not in git they will not be online.

The first push carries roughly 53 MB and will take a few minutes
(Section~\ref{sec:size}). Later pushes only carry what changed, so they are
quick.

\subsection{Visitor analytics}\label{sec:analytics}

Counting is done by GoatCounter; the dashboard is
\url{https://alignain.goatcounter.com}. It is privacy-preserving, needs no
cookie banner, and the free tier is ample for a personal site.

\paragraph{Where the tag is defined.} Once, in \file{\_quarto.yml} under
\texttt{format > html > include-in-header}. Nothing else should ever hard-code
the endpoint --- see the hook below, which reads it back out of that file
precisely so the two cannot drift apart.

\paragraph{Why a post-render hook exists.} \texttt{include-in-header} only
reaches pages that \emph{Quarto renders} --- nine of them. The rest of the
site is finished HTML that Quarto merely \emph{copies}: the fourteen
\file{Rclass\_2025} lecture decks, the hand-rendered inflation tracker, the
Shinylive splicer app. Those were being served untracked. \file{R/post-render.R}
walks \file{\_site} after every render and splices the tag into any page that
lacks it:

\begin{center}
\begin{tabular}{@{}lr@{}}
\toprule
\textbf{\texttt{post-render:} in \file{\_quarto.yml}} & \textbf{Pages counted} \\
\midrule
commented out & 10 of 39 \\
enabled       & \textbf{39 of 39} \\
\bottomrule
\end{tabular}
\end{center}

It is safe to re-run: pages that already carry the tag are skipped, so
nothing is ever injected twice. It works on raw bytes rather than
\texttt{readLines()}/\texttt{writeLines()} because the lecture decks are not
all UTF-8 and a read-then-write round trip would re-encode them.

\paragraph{The footer badge.} The counter in the page footer is written as a
raw \texttt{<img>} in \file{\_quarto.yml}, \emph{not} as a markdown image:
\begin{lstlisting}
right: '<img class="visitor-badge" src="...TOTAL.svg?no_branding=1"
             alt="Visitor count">'
\end{lstlisting}
Written the markdown way, \texttt{!\lbrack alt\rbrack(url)}, Quarto wraps it in
a \texttt{<figure>} and prints the alt text as a visible caption --- so the
footer showed the words ``visitor count'' underneath a badge that already
read ``Views for this site''. As raw HTML the text survives in \texttt{alt}
for screen readers without being drawn twice.

The badge SVG hard-codes purple (\texttt{\#9a15a4}) for its border and text
and ignores a \texttt{?style=} parameter --- the server returns the same
colours regardless --- so it is toned down from our side, in
\file{styles.css} under \texttt{.visitor-badge}: \texttt{grayscale(1)},
\texttt{opacity 0.5} and a 38\,px height, lifting to \texttt{0.85} on hover
for anyone who actually wants to read the number.

\paragraph{Two things that look like faults and are not.}
\begin{itemize}[itemsep=2pt]
\item \textbf{The badge reads 0 locally.} \file{count.js} deliberately
  ignores \texttt{localhost} and \texttt{file://}, so no local preview has
  ever registered a visit. The number only becomes real once the site is
  hosted.
\item \textbf{The badge was a broken image.} The \texttt{/counter/} endpoint
  returns HTTP 403 until \emph{Settings $\to$ allow public access to the
  counter} is ticked in GoatCounter. That has been enabled; if the badge ever
  reverts to a broken-image icon, check that setting first.
\end{itemize}

\section{Troubleshooting}

\begin{description}[style=nextline,itemsep=4pt]
\item[Stats badges show NA or old numbers]
  Google blocked the live fetch and no cache exists. Update the
  \emph{scholar\_fallback} sheet, or render again later (a successful fetch
  writes the cache and fixes it permanently).
\item[\texttt{PermissionDenied \ldots \_site} during render]
  Something (antivirus, sync client, an open file) is locking
  \file{\_site}. Close it, delete \file{\_site}, render again. This is
  chronic inside Google Drive folders --- keep the project out of Drive.
\item[\texttt{Unable to locate an installed version of R}]
  Set \texttt{QUARTO\_R} as in Section~\ref{sec:render}, or render from
  RStudio.
\item[A page renders empty, and the log says \texttt{Sheet '...' is missing}]
  A sheet was deleted or renamed in \file{cv\_data.xlsx}. \file{R/helpers.R}
  deliberately warns and renders that section empty rather than killing the
  whole render, so \textbf{watch the render log} --- the page will look fine
  apart from the missing content. Re-add the sheet with its original name
  and column headings. (This is how the \emph{software} sheet went missing
  in July 2026, emptying the Software page --- restored in August 2026.
  Every rewrite of the workbook leaves a dated backup beside it in
  \file{cvnain/data/}, which is what to fall back on.)
\item[A CSS rule has no effect at all]
  First suspect the selector, not the CSS. Quarto renames listing
  containers (\texttt{\#course-pages} becomes
  \texttt{\#listing-course-pages}), and a rule aimed at the wrong id fails
  \emph{silently} --- nothing errors, the page just renders unstyled. Open
  the rendered file in \file{\_site/} and search for the class or id you
  are targeting before assuming the rule is wrong.
\item[A slide link works locally but 404s on the live site]
  Almost always capitalisation --- see Section~\ref{sec:slidelinks}.
  Windows ignores it; the GitHub Pages server does not.
\item[The Action is green but the site is 404]
  \emph{Settings} $\to$ \emph{Pages} $\to$ \emph{Source} is still on
  \emph{Deploy from a branch}. It must be \textbf{GitHub Actions}. Nothing
  is published until it is.
\item[The deploy step fails with a permissions error]
  Same cause as above --- \texttt{actions/deploy-pages} cannot create a
  deployment while Pages is set to branch mode.
\item[\texttt{git push} rejected: \emph{updates were rejected}]
  The GitHub repository was created with a README or licence, so it has a
  commit the local folder does not. Easiest fix is to delete the repository
  and recreate it empty; alternatively \texttt{git pull --rebase origin main}
  first.
\item[The homepage surface looks smooth and invented rather than estimated]
  \file{irf-surface.js} did not load. Check it is committed, that
  \file{\_quarto.yml} still lists it under \texttt{resources}, and that
  \file{index.qmd} loads it \emph{before} \file{background.js}.
\item[\texttt{git push} rejected: file exceeds 100 MB]
  Nothing in the project is near that any more (Section~\ref{sec:size}).
  If it happens, something generated has been committed by mistake ---
  check that \file{software/splicer-app/} is still ignored.
\item[The splicer app is live but GoatCounter shows no visits for it]
  \file{shinylive::export()} writes its own \file{index.html}, so the
  analytics tag has to be re-applied after every export.
  \file{software/build-shinylive.R} does this automatically, and
  \file{R/post-render.R} catches it again at the end of every render
  (Section~\ref{sec:analytics}); if the page was rebuilt some other way,
  re-run that script.
\item[Lecture decks or the tracker are not being counted]
  \texttt{post-render:} is commented out in \file{\_quarto.yml}. Coverage
  drops from 39 pages to 10 --- the copied pages are the ones that lose the
  tag (Section~\ref{sec:analytics}). The render log prints the count every
  time: \texttt{post-render: analytics tag --- \ldots\ added}.
\item[Cloud build fails on a missing R package]
  Add the package name to the \texttt{install.packages(...)} line in
  \file{.github/workflows/publish.yml}. It currently installs
  \texttt{readxl, dplyr, glue, plotly, httr2, jsonlite, rvest, httr,
  tidyr, knitr, rmarkdown, shinylive}.
\item[Data Corner chart is stale or empty]
  The render log says \texttt{MoSPI: live fetch failed \ldots\ using cache}.
  The API was unreachable and \file{data/mospi\_cache.csv} was drawn
  instead; if that file is missing too, the chart cannot render. Re-render
  when the API is back --- one success repopulates the cache
  (Section~\ref{sec:mospi}).
\item[CV renders but the site still serves an old \file{cv.pdf}]
  The site copies \file{cvnain/Dr-Nain\_CV.pdf} at render time
  (\file{R/pre-render.R}); re-render the site after re-rendering the CV.
\end{description}

\end{document}
